% Part: normal-modal-logic
% Chapter: sequent-calculus
% Section: introduction

\documentclass[../../../include/open-logic-section]{subfiles}

\begin{document}

\olfileid[zh]{nml}{seq}{int}

\olsection{引言}

命题逻辑的矢列演算可以通过处理~\iftag{prvBox}{$\Box$\iftag{prvDiamond}{ 与~}{}}{}%
\iftag{prvDiamond}{$\Diamond$}{} 的附加规则来扩充。例如，对~$\Log{K}$，我们
有 \Log{LK} 再加上：
\[
  \iftag{prvBox}{\iftag{prvDiamond}{% ◇ 与 □ 都是原始符号
    \Axiom$\Gamma \fCenter \Delta, !A$
    \RightLabel{$\Box$}
    \UnaryInf$\Box\Gamma \fCenter \Diamond\Delta, \Box!A$
    \DisplayProof
    \qquad
    \Axiom$!A, \Gamma \fCenter \Delta$
    \RightLabel{$\Diamond$}
    \UnaryInf$\Diamond!A, \Box\Gamma \fCenter \Diamond\Delta$
    \DisplayProof
}{% 只有 □ 是原始符号
\Axiom$\Gamma \fCenter !A$
\RightLabel{$\Box$}
\UnaryInf$\Box\Gamma \fCenter \Box!A$
\DisplayProof
}}{% 只有 ◇ 是原始符号
  \Axiom$!A \fCenter \Delta$
  \RightLabel{$\Diamond$}
  \UnaryInf$\Diamond!A \fCenter \Diamond\Delta$
  \DisplayProof
}
\]
对~$\Log{K}$ 的扩张，还必须再添加附加规则。

并非每种模态逻辑都有这样的矢列演算。甚至~$\Log{S5}$——它在语义上很简单（它完全可以不用可及关系来定义）——也还不知道是否有由~$\Log{LK}$ 得到、在无~\Cut 规则下完全的矢列演算。不过，它有一个无割且完全的 \emph{超矢列}演算。

\end{document}
